\section{Chương~\ref{sec:neurontypes}. Các loại nơ-ron}

Các con số của chương dựa trên những phép đếm tế bào trực tiếp trong não người ở những nơi có phép đếm như vậy; quy tắc ``một nơ-ron --- một chất dẫn truyền'' dựa trên các bản đồ tế bào và những thí nghiệm đã tìm ra cả ngoại lệ; vai trò của \nt{DOPA} dựa trên các bản ghi xung; còn vai trò của các kênh quảng bá khác dựa trên giả thuyết.

{\footnotesize
\setlength{\tabcolsep}{3pt}\renewcommand{\arraystretch}{1.15}
\begin{longtable}{>{\raggedright}p{0.5cm}>{\raggedright}p{4.0cm}>{\raggedright}p{5.6cm}>{\raggedright}p{2.3cm}>{\raggedright\arraybackslash}p{1.7cm}}
\toprule
TT & Khẳng định & Thí nghiệm và điều đã đo & Nguồn & Trạng thái \\
\midrule\endfirsthead
\toprule
TT & Khẳng định & Thí nghiệm và điều đã đo & Nguồn & Trạng thái \\
\midrule\endhead
1 & Não người có khoảng $8.6\cdot10^{10}$ nơ-ron, và gần như mọi nơ-ron \nt{GLUT} là nơ-ron nhỏ của tiểu não (bảng~\ref{tab:types}) & Phương pháp phân đoạn đẳng hướng: mô được hòa tan thành huyền phù nhân đồng nhất, rồi đếm các nhân mang dấu ấn nơ-ron NeuN. Ở nam giới trưởng thành có $86.1\pm8.1$ tỷ nơ-ron; trong vỏ đại não chỉ có $19\,\%$, phần lớn nằm ở tiểu não & \cite{Azevedo2009} & đã đo \\
2 & Hầu như mỗi nơ-ron có một chất dẫn truyền chính (mệnh đề~\ref{prop:dale}), nhưng có ngoại lệ & Bản đồ phiên mã não chuột: khoảng $4\cdot10^6$ tế bào, $5322$ cụm; mỗi cụm được gán chất dẫn truyền theo các gen tổng hợp và đóng gói. Ngoại lệ được đo trực tiếp: trong lát cắt, nơ-ron \nt{DOPA} nhả cả GABA qua cùng chất vận chuyển túi và ức chế nơ-ron thể vân; ở một phần sợi trục \nt{DOPA}, glutamate và dopamine đi ra từ những đầu tận cùng khác nhau của cùng một dây (kính hiển vi điện tử và quang di truyền) & \cite{Strata1999,Yao2023,Tritsch2012,Zhang2015} & đã đo \\
3 & Cả cột ma trận trọng số cùng một dấu~\eqref{eq:dalesign} & Suy luận trong chương từ mệnh đề~\ref{prop:dale} và từ việc gần như mọi thụ thể glutamate là kích thích, còn thụ thể GABA là ức chế. Chưa có kiểm chứng trực tiếp ``mọi nơi nhận của một nơ-ron đều nhận trọng số cùng dấu'' như một quy tắc chung; phản ví dụ là việc nơ-ron \nt{DOPA} nhả kèm GABA (dòng~2) và những nơi nhận có thụ thể trái dấu (dòng~11) & suy luận trong chương & lý thuyết \\
4 & Từ ba phần tư đến bốn phần năm nơ-ron vỏ não là \nt{GLUT}, từ một phần năm đến một phần tư là \nt{GABA} & Nhuộm miễn dịch GABA trong các cột rộng $50\um$ xuyên suốt độ dày vỏ não ở $10$ vùng vỏ não khỉ: ở $8$ vùng, nơ-ron GABA chiếm khoảng $25\,\%$ tổng số nơ-ron, ở vỏ thị giác sơ cấp --- dưới $20\,\%$ & \cite{Hendry1987} & đã đo \\
5 & Nơ-ron \nt{GLUT} có khoảng $10^4$ đầu ra & Lập thể học trên tử thi: vỏ não mới của nam thanh niên có $1.64\cdot10^{14}$ khớp thần kinh (kính hiển vi điện tử, disector) và khoảng $2.3\cdot10^{10}$ nơ-ron. Tỷ số cho $\sim7\cdot10^3$ khớp thần kinh trên một nơ-ron; tính trung bình, số đầu vào bằng số đầu ra & \cite{Tang2001synapses,Pakkenberg1997} & đã đo \\
6 & Đầu ra của nơ-ron \nt{DOPA} phân nhánh thành $10^5$--$10^6$ đầu tận cùng; đây là ước tính theo độ phủ vùng đích chứ không phải phép đếm & Đánh dấu màng bằng virus ở từng nơ-ron \nt{DOPA} của chuột cống và tái dựng toàn bộ sợi trục: một nơ-ron phủ từ $0.45$ đến $5.7\,\%$ (trung bình $2.7\,\%$) thể tích thể vân. Đại lượng đo được là độ phủ chứ không phải số đầu tận cùng: số này được suy từ độ phủ theo bậc độ lớn, chưa có phép đếm trực tiếp. Với \nt{SERO} và \nt{NORA}, số đầu tận cùng là ước tính thô & \cite{Matsuda2009} & đã đo \\
7 & Nơ-ron quảng bá rất ít: \nt{DOPA} $\sim5\cdot10^5$ (nhóm chính trong hai nhóm, cận dưới), \nt{SERO} $\sim3\cdot10^5$ (tổng của hai phép đếm từng phần), \nt{HIST} $\sim6\cdot10^4$ (bảng~\ref{tab:types}, hình~\ref{fig:typepie}) & Các phép đếm lập thể học ở người: $550\,000$ nơ-ron sắc tố trong chất đen (không tính vùng trần não thất); $235\,000$ nơ-ron trong nhân đường giữa lưng (mọi nơ-ron của nhân) và thêm $125\,000$ nơ-ron serotonin ở trần cầu não; khoảng $64\,000$ nơ-ron histamine ở vùng dưới đồi. Với \nt{NORA}, \nt{GLYC}, \nt{ACET} và \nt{ADRE} không có phép đếm trực tiếp: trong bảng đó là ước tính thô theo bậc độ lớn & \cite{Pakkenberg1991,Baker1990,Baker1991,Airaksinen1991} & đã đo \\
8 & Glutamate trong khe vọt lên khoảng một milimol và giảm trong $\sim1\ms$ & Qua sự đẩy một chất chẹn cạnh tranh tách ra nhanh khỏi thụ thể NMDA trong lúc truyền qua khớp thần kinh (mẫu nuôi cấy nơ-ron hồi hải mã): đỉnh $1.1\mM$, hằng số suy giảm $1.2\ms$ & \cite{Clements1992} & đã đo \\
9 & Trong vỏ não đang làm việc, dòng kích thích và dòng ức chế gần như cân bằng, nơ-ron nằm sát ngưỡng & Lý thuyết: mạng có liên kết mạnh tự đi vào chế độ cân bằng. Thực nghiệm: ở vỏ não trước trán của chồn sương in vivo, trong các trạng thái ``lên'', điện thế đảo chiều của dòng qua khớp thần kinh giữ quanh $-37\mV$ suốt hàng trăm mili giây, dù độ dẫn thay đổi $\sim21$\,nS: kích thích và ức chế cùng tăng và cùng giảm & \cite{vanVreeswijk1996,Haider2006} & đã đo \\
10 & Nơ-ron \nt{DOPA} báo sai số dự đoán phần thưởng~\eqref{eq:rpe} (hình~\ref{fig:rpe}) & Xung của nơ-ron \nt{DOPA} ở khỉ: chùm xung với nước ép bất ngờ; sau khi học --- chùm xung với tín hiệu báo trước và không đáp ứng với nước ép đã được dự đoán; tần số sụt khi nước ép đã hứa không được cho. Trong thí nghiệm định lượng, tần số tỷ lệ với hiệu giữa phần thưởng hiện tại và trung bình có trọng số của các phần thưởng trước, nhưng chỉ với sai số dương. Bản thân phương trình~\eqref{eq:rpe} là thuật toán sai phân thời gian & \cite{Schultz1997,Bayer2005,SuttonBarto2018} & đã đo \\
11 & Dấu và tốc độ do thụ thể quyết định: thụ thể hướng ion --- phần nhỏ của mili giây, thụ thể hướng chuyển hóa --- chậm hơn nhiều (mệnh đề~\ref{prop:receptor}) & Xung glutamate ngắn lên các mảnh màng của nơ-ron hồi hải mã: dòng qua thụ thể hướng ion tăng (từ $20$ đến $80\,\%$) trong $0.2$--$0.6\ms$ và giảm trong $2$--$3\ms$. Điện thế ức chế chậm của nơ-ron tháp bị xóa bởi chất chẹn chọn lọc thụ thể GABA hướng chuyển hóa. Hai họ thụ thể dopamine có tác động ngược nhau; ở thể vân, nơ-ron có thụ thể D1 cần dopamine để tăng cường khớp thần kinh, nơ-ron có D2 --- để làm yếu khớp thần kinh & \cite{Colquhoun1992,DutarNicoll1988,Beaulieu2011,Shen2008} & đã đo \\
12 & Kênh quảng bá không phải một dây dẫn mà là một bó ít đường (mục~\ref{sec:buses}) & Đánh dấu di truyền bằng virus các nơ-ron serotonin của nhân đường giữa lưng ở chuột nhắt: nơ-ron gửi đầu ra tới hạch hạnh nhân và tới vỏ não trán nằm ở những chỗ khác nhau trong nhân, phân nhánh vào những vùng bổ sung cho nhau và đáp ứng ngược chiều với kích thích khó chịu. Tổng quan: các phân nhóm nơ-ron của nhân lục (\nt{NORA}) mã hóa những quá trình khác nhau & \cite{Ren2018,Poe2020} & đã đo \\
13 & \nt{NORA} đặt hệ số khuếch đại $g$ & Giả thuyết khuếch đại thích nghi; mô hình mạng trong đó catecholamine làm tăng độ dốc của đặc tuyến truyền đạt. Chưa có phép đo trực tiếp ``$g$ tăng theo noradrenaline'' trên toàn mạng; đã đo được rằng đường kính đồng tử bám theo xung của nơ-ron nhân lục ở khỉ & \cite{AstonJones2005,ServanSchreiber1990,Joshi2016} & giả thuyết \\
14 & \nt{ACET} chuyển ``ghi/đọc'' và đặt tốc độ học & Giả thuyết. Cơ sở: trong lát cắt vỏ khứu giác chuột cống, chất chủ vận acetylcholine ($100\,\mu\text{M}$) làm yếu các khớp thần kinh bên trong, ``hồi tưởng'', hơn $60\,\%$, còn các khớp thần kinh đầu vào --- dưới $15\,\%$ & \cite{Hasselmo2006,HasselmoBower1992} & giả thuyết \\
15 & \nt{SERO} đặt hệ số chiết khấu $\gamma$; nơ-ron serotonin làm việc đều, vài xung mỗi giây, và gần như im lặng trong một pha của giấc ngủ & Giả thuyết ``serotonin $=\gamma$''. Lập luận mạnh nhất: hoạt hóa quang di truyền các nơ-ron serotonin của nhân đường giữa lưng ở chuột nhắt làm giảm số lần bỏ chờ sớm phần thưởng trì hoãn. Tần số xung và sự im lặng trong giấc ngủ có cử động mắt nhanh đã được đo (tổng quan các bản ghi) & \cite{Doya2002,Miyazaki2014,JacobsAzmitia1992} & giả thuyết \\
\bottomrule
\end{longtable}}
